% ----------------------------------------------------------------------------
\chapter{Sobolev, Poincar\'e and relative isoperimetry}\label{ch:sobolev}
\module{NoCompromise.Sobolev.BV,
        NoCompromise.Sobolev.Extension,
        NoCompromise.Sobolev.Maximal,
        NoCompromise.Sobolev.Poincare,
        NoCompromise.Sobolev.RelativeCubes,
        NoCompromise.Sobolev.RelativeScaling}

\section{The BV Sobolev inequality}

\begin{lemma}[Gagliardo--Nirenberg in $\R^3$, smooth case]\label{lem:GN-smooth}
For $f\in C^1_c(\R^3)$,
\begin{equation}\label{eq:GN-smooth}
  \|f\|_{L^{3/2}(\R^3)}\le C\prod_{i=1}^3\|\partial_if\|_{L^1}^{1/3}
  \le \tfrac{C}{3}\sum_{i=1}^3\|\partial_if\|_{L^1}.
\end{equation}
\end{lemma}

\begin{proof}
Put $A_i(x_1,\dots,\widehat{x_i},\dots,x_3):=\int_\R|\partial_if|\,dx_i$.  The
fundamental theorem of calculus gives $|f(x)|\le A_i$ for each $i$, hence
$|f|^{3/2}\le(A_1A_2A_3)^{1/2}$.  Fubini and Cauchy--Schwarz successively in
$x_1,x_2,x_3$ give
$\int_{\R^3}(A_1A_2A_3)^{1/2}\le\prod_i\|\partial_if\|_{L^1}^{1/2}$.  Raise to
the power $2/3$ and use the arithmetic--geometric mean inequality.
\end{proof}

\begin{theorem}[BV Sobolev inequality]\label{thm:bv-sobolev}
\uses{def:variation}
For $f\in BV(\R^3)$,
\begin{equation}\label{eq:bv-sobolev}
  \|f\|_{L^{3/2}(\R^3)}\le C\,|Df|(\R^3),
\end{equation}
with $C$ absolute.  In particular, for $E$ of finite measure and finite
perimeter, $|E|^{2/3}\le C\Per(E)$.
\end{theorem}

\begin{proof}\uses{lem:GN-smooth,lem:strict-approx}
\eqref{eq:GN-smooth} plus strict approximation
(Theorem~\ref{lem:strict-approx}) and Fatou.
\end{proof}

\begin{lemma}[The planar analogue]\label{lem:planar-sobolev-L2}
For $f\in C^1_c(\R^2)$,
\begin{equation}\label{eq:planar-sobolev-L2}
  \|f\|_{L^2(\R^2)}^2\le\|\partial_1f\|_{L^1}\|\partial_2f\|_{L^1},
  \qquad
  \|f\|_{L^2(\R^2)}\le\tfrac12\sum_{i=1}^2\|\partial_if\|_{L^1},
\end{equation}
For $f\in BV(\R^2)$ the precise measure version is
\[
  \|f\|_{L^2}^2\le |D_1f|(\R^2)|D_2f|(\R^2),\qquad
  \|f\|_{L^2}\le\tfrac12\bigl(|D_1f|(\R^2)+|D_2f|(\R^2)\bigr)
  \le |Df|(\R^2).
\]
\end{lemma}

\begin{proof}\uses{lem:strict-approx,lem:GN-smooth}
The identical two-coordinate fundamental-theorem-of-calculus calculation as
in Lemma~\ref{lem:GN-smooth} proves the two inequalities for
$C^1_c(\R^2)$.  For $f\in BV(\R^2)$ put $h_\varepsilon:=\rho_\varepsilon*f$.
The distributional convolution identities give
\[
  \|\partial_i h_\varepsilon\|_{L^1}
  \le |D_if|(\R^2),\qquad i=1,2.
\]
Apply the smooth inequalities first to $\chi_Rh_\varepsilon$ and let
$R\to\infty$; the cutoff error is at most
$C R^{-1}\|h_\varepsilon\|_1$ and hence vanishes.  Finally
$h_\varepsilon\to f$ in $L^1$, so a subsequence converges a.e.; Fatou gives
the claimed coordinatewise measure bounds and then the arithmetic--geometric
mean bound.  This coordinatewise argument is stronger than isotropic strict
convergence and does not infer convergence of $|D_if|$ from
Theorem~\ref{lem:strict-approx}.
\end{proof}

\section{Extension and Poincar\'e}

\begin{lemma}[Bi-Lipschitz chain rule in BV]\label{lem:bilip-chain}
\uses{def:variation}
Let $U,V\subset\R^n$ be open, let $\Phi:U\to V$ be a bi-Lipschitz
homeomorphism, let $D\subset V$ be open, and let $g\in BV(D)$.  Then
$g\circ\Phi\in BV(\Phi^{-1}(D))$ and
\[
 \|g\circ\Phi\|_{L^1(\Phi^{-1}(D))}
 +|D(g\circ\Phi)|(\Phi^{-1}(D))
 \le C\bigl(\|g\|_{L^1(D)}+|Dg|(D)\bigr),
\]
where $C$ depends only on $n$, $\Lip(\Phi)$, and $\Lip(\Phi^{-1})$.
\end{lemma}

\begin{proof}\uses{lem:strict-approx,lem:lsc}
Take a strict smooth approximation $g_j$ of $g$.  Ordinary change of
variables and
$\nabla(g_j\circ\Phi)=D\Phi^{\mathsf T}(\nabla g_j)\circ\Phi$ bound both the
$L^1$ norm and the variation of $g_j\circ\Phi$; lower semicontinuity
(Lemma~\ref{lem:lsc}) passes the estimate to $g\circ\Phi$.
\end{proof}

\begin{lemma}[BV extension]\label{lem:bv-extension}
\uses{def:variation,lem:bilip-chain,lem:bv-product-smooth}
Let $D\subset\R^n$ ($n=2,3$) be a bounded connected Lipschitz domain.  There
is a bounded linear $\mathcal E:BV(D)\to BV(\R^n)$, with image supported in
one fixed bounded neighbourhood of $\overline D$, such that
\begin{equation}\label{eq:bv-extension}
  \|\mathcal Eg\|_{L^1(\R^n)}+|D\mathcal Eg|(\R^n)
  \le C_D\bigl(\|g\|_{L^1(D)}+|Dg|(D)\bigr).
\end{equation}
\end{lemma}

\begin{proof}\uses{lem:bilip-chain,lem:bv-product-smooth}
First consider a flat half-cube
$Q^+:=(-1,1)^{n-1}\times(0,1)$ and a smaller cube $Q_0\Subset(-1,1)^n$.
For $g\in BV(Q^+)$ define its even reflection by
$Rg(x',x_n):=g(x',|x_n|)$.  If
$X\in C_c^1(Q_0;\R^n)$, split
$\int Rg\,\dvg X$ over the two half-cubes and change $x_n$ to $-x_n$ in the
lower one.  The result is the pairing of $g$ with the divergence of the
upper test field plus the reflected lower test field, with the normal
component reflected with a minus sign.  Its pointwise norm is at most
$2\|X\|_\infty$.  Definition~\ref{def:variation} therefore gives
\[
 \|Rg\|_{L^1(Q_0)}\le2\|g\|_{L^1(Q^+)},\qquad
 |D(Rg)|(Q_0)\le2|Dg|(Q^+).
\]
This direct calculation is the reason no trace theorem for Lipschitz domains
is being assumed.

Cover $\partial D$ by finitely many Lipschitz graph charts, flatten each by a
bi-Lipschitz map, use the preceding even reflection, multiply by a smooth
partition of unity, and extend the piece supported strictly inside $D$ by
zero.  Lemma~\ref{lem:bilip-chain} controls each flattening,
Lemma~\ref{lem:bv-product-smooth} controls the cutoffs, and the finite sum of
the extended pieces defines the required bounded linear map.  Choosing all
cutoffs in one fixed bounded collar gives the support assertion.
\end{proof}

\begin{theorem}[BV--Poincar\'e]\label{thm:bv-poincare}
\uses{lem:bv-extension}
Let $D$ be a bounded connected Lipschitz domain in $\R^n$, $n=2,3$, and
$g_D:=\avg_Dg$.  Then
\begin{equation}\label{eq:bv-poincare}
  \|g-g_D\|_{L^1(D)}\le C_D\,|Dg|(D),
  \qquad
  \|g-g_D\|_{L^{n/(n-1)}(D)}\le C_D\,|Dg|(D).
\end{equation}
\end{theorem}

\begin{proof}
\uses{lem:bv-extension,thm:bv-compactness,thm:bv-sobolev,lem:planar-sobolev-L2}
Suppose the first inequality failed: there would be $f_j\in BV(D)$ with
$(f_j)_D=0$, $\|f_j\|_{L^1(D)}=1$ and $|Df_j|(D)\to0$.  By
\eqref{eq:bv-extension} the $\mathcal Ef_j$ are bounded in $BV(\R^n)$ and
supported in one fixed compact set; apply the grid compactness argument of
Theorem~\ref{thm:bv-compactness} on a cube containing that set.  A
subsequence converges in $L^1(\R^n)$, hence in $L^1(D)$, to some $f$.  Testing
distributional derivatives compactly inside $D$ gives $Df=0$, so $f$ is
constant on balls and hence on connected $D$; mean zero makes the constant
zero, contradicting $\|f\|_{L^1(D)}=1$.  The second inequality follows by
applying Theorem~\ref{thm:bv-sobolev} (or
Lemma~\ref{lem:planar-sobolev-L2} when $n=2$) to
$\mathcal E(g-g_D)$ and using the first.
\end{proof}

\section{Relative isoperimetric inequalities}

\begin{remark}[Three distinct signatures]\label{rem:three-signatures}
The cube, ball and planar-disk cases below are \emph{separate statements} with
separate constants.  The ball case is \emph{not} obtained from the cube case
by inclusion --- the relative perimeter on the right would be the wrong one.
Two of the three are used: \eqref{eq:rel-iso-cube} in
Theorem~\ref{thm:compactness}, and \eqref{eq:rel-iso-ball} in
Lemma~\ref{lem:ahlfors} and Lemma~\ref{lem:open-representative}.
\eqref{eq:rel-iso-disk} is \emph{not} on the critical path; it is recorded for
completeness of the planar toolkit and may be skipped
(Formalisation note~\ref{fnote:off-critical-path}).
\end{remark}

\begin{proposition}[Relative isoperimetry, cube]\label{prop:rel-iso-cube}
\uses{thm:bv-poincare}
For a cube $Q\subset\R^3$ of side $\ell$ and a measurable $E$ with
$\ind E\in BV(Q)$,
\begin{equation}\label{eq:rel-iso-cube}
  \min\{|E\cap Q|,|Q\setminus E|\}^{2/3}\le C_{\mathrm{cube}}\,\Per(E;Q),
\end{equation}
with $C_{\mathrm{cube}}$ absolute (scale-invariant).
\end{proposition}

\begin{proof}\uses{thm:bv-poincare}
Apply \eqref{eq:bv-poincare} on $Q$ to $g=\ind E$, and put
$\theta:=|E\cap Q|/|Q|$.  If $\theta\le1/2$, then
$|\ind E-\theta|\ge1/2$ on $E\cap Q$; if $\theta\ge1/2$, apply the same
argument to $Q\setminus E$.  Scaling the unit-cube inequality gives
\eqref{eq:rel-iso-cube} with an absolute constant.
\end{proof}

\begin{proposition}[Relative isoperimetry, ball]\label{prop:rel-iso-ball}
\uses{thm:bv-poincare}
For a ball $B\subset\R^3$ and a measurable $E$ with $\ind E\in BV(B)$,
\begin{equation}\label{eq:rel-iso-ball}
  \min\{|E\cap B|,|B\setminus E|\}^{2/3}\le C_{\mathrm{ball}}\,\Per(E;B),
\end{equation}
with $C_{\mathrm{ball}}$ absolute.
\end{proposition}

\begin{proof}\uses{thm:bv-poincare}
Apply \eqref{eq:bv-poincare} on $B$ to $g=\ind E$ and use the same two cases
$|E\cap B|\le|B|/2$ and $|B\setminus E|\le|B|/2$.  Dilation from the unit
ball gives \eqref{eq:rel-iso-ball} with an absolute constant.
\end{proof}

\begin{proposition}[Relative isoperimetry, planar disk]\label{prop:rel-iso-disk}
\uses{thm:bv-poincare}
For a planar disk $D\subset\R^2$ and a measurable $A\subset\R^2$ with
$\ind A\in BV(D)$,
\begin{equation}\label{eq:rel-iso-disk}
  \min\{|A\cap D|,|D\setminus A|\}^{1/2}\le C\,\Per_2(A;D).
\end{equation}
\end{proposition}

\begin{proof}\uses{thm:bv-poincare}
Apply \eqref{eq:bv-poincare} in dimension two on $D$ to $g=\ind A$.  Splitting
according as $|A\cap D|\le|D|/2$ or $|D\setminus A|\le|D|/2$, and then
scaling the unit-disk inequality, gives \eqref{eq:rel-iso-disk}.
\end{proof}

\section{The weak maximal estimate}\label{sec:maximal}

\begin{definition}[Planar maximal density]\label{def:maximal}
For a finite Radon measure $\mu$ on $\R^2$ put
$M\mu(x):=\sup_{r>0}\mu(B_r^2(x))/(\pi r^2)$, and for $R>0$ put
$M_R\mu(x):=\sup_{0<r<R}\mu(B_r^2(x))/(\pi r^2)$.
\end{definition}

\begin{lemma}[Weak maximal estimate]\label{lem:maximal}
\uses{def:maximal}
For a finite Radon measure $\mu$ on $\R^2$ and $t>0$,
\begin{equation}\label{eq:maximal}
  \bigl|\{M\mu>t\}\bigr|\le\frac{25}{t}\,\mu(\R^2),
\end{equation}
and the same bound holds for $M_R\mu$.
\end{lemma}

\begin{proof}
Let $K\subset\{M\mu>t\}$ be compact and choose for each $x\in K$ a ball
$B_x=B_{r_x}^2(x)$ with $\mu(B_x)>t\pi r_x^2$.  The five-covering lemma
(a consequence of \ref{base:besicovitch}) gives a disjoint subfamily $B_i$
whose fivefold dilates cover $K$, so
$|K|\le25\sum_i\pi r_i^2<\frac{25}{t}\sum_i\mu(B_i)\le\frac{25}{t}\mu(\R^2)$.
Inner regularity removes the compactness assumption.  Restricting radii gives
the $M_R$ version.
\end{proof}

